% Section 10.1, from lectures/L10/appendix/a_bench.md.

\section{The Two-Link Bench}
\label{c10:sec:bench}\label{c10:app:a}
Everything the book built meets on one bench: your VHDL peripheral on the DE0-CV, your C++ driver on
the Nano's ATmega328P, and a PC terminal at the far end. What makes the bench worth understanding is
that \textbf{two different serial links} run across it, for two different jobs, and confusing them is
the most common bring-up mistake.

\textbf{The control plane is SPI.} The ATmega328P is the CPU, and it configures and drives the
peripheral by reading and writing its registers over SPI, in the five-byte transactions of Part~3 of
the protocol (\secref{proto:part3}). This is \code{AvrSpi} on one end and the provided
\code{spi_slave} and \code{spi_reg_bridge} on the other, and it carries \emph{register} traffic, not
user data.

\textbf{The data plane is the UART itself.} The peripheral's own \code{tx} and \code{rx} pins carry
bytes to and from the outside world, out to a USB-serial adapter and a PC terminal. This is the line
the whole peripheral exists to drive.

So a byte the user types travels \textbf{in} on the data plane (UART \code{rx}), is read by the CPU
\textbf{over the control plane} (SPI, \code{RX_DATA}), and, in the echo application, is written back
\textbf{over the control plane} (SPI, \code{TX_DATA}) to leave again on the data plane (UART
\code{tx}). Two protocols, one byte.

\subsection{Level shifting: 5\,V meets 3.3\,V}
\label{c10:sec:level_shifting}
The Nano runs at \textbf{5\,V}, the DE0-CV at \textbf{3.3\,V}, and that mismatch is the one
electrical hazard on the bench.

The \textbf{four SPI lines cross a level shifter}, a bidirectional BSS138-style board: \code{SCK},
\code{MOSI} and \code{SS} go from 5\,V on the Nano down to 3.3\,V at the FPGA, while \code{MISO} goes
from 3.3\,V up to 5\,V. The shifter's HV side ties to the Nano's 5\,V, its LV side to the DE0-CV's
3.3\,V, and the grounds are common. The \textbf{data-plane UART needs no shifting} at all, because
the DE0-CV's \code{tx} and \code{rx} are already 3.3\,V and the USB-serial adapter is set to 3.3\,V
logic, so those two connect directly.

Drive a 3.3\,V FPGA input straight from a 5\,V pin and, best case, it misbehaves; worst case, you
damage the pin. The ladder's control-plane rung is where a crossed or unshifted SPI line first shows
itself, but a rung can only report the damage, not prevent it: by the time any rung runs, the line
has been powered. That is why the wiring is checked before power is applied
(Exercise~\ref{c10:ex:c1}).

\subsection{The board wrapper and Quartus}
\label{c10:sec:wrapper}
\code{uart_top} passed its system testbench in simulation (\chapref{5}), but a testbench is not a pin
assignment. To run on real silicon, \code{uart_top} is wrapped by \textbf{\file{uart_board.vhd}}, the
Quartus top level that maps the peripheral's ports (\code{clock}, \code{reset_n}, the SPI lines,
\code{rx}, \code{tx}) onto specific DE0-CV pins: an on-board 50\,MHz clock circuit, a reset button,
and GPIO-header pins for SPI and the UART.

That wrapper is board I/O, not peripheral logic, so it lives with the Quartus project rather than in
\file{hw/}, and it is provided. You already synthesized it and programmed the board in \chapref{9},
so nothing about the FPGA toolchain is new here; the only change now is the loopback edit each rung
of the ladder asks for.

\subsection{Why a ladder, and what integration proves}
\label{c10:sec:ladder}
The peripheral works in simulation and the driver passes its host tests, yet the assembled bench can
still fail, because simulation and host tests cannot see a level shifter, a crossed wire, a clock
that is not really 50\,MHz, or SPI and UART timing on real edges. Integration is the top of the test
pyramid: it is the only level that exercises the actual chip boundary.

The \textbf{bring-up ladder} climbs that gap one rung at a time, each rung proving something the
last could not, so that a failure points at one layer instead of the whole bench. \textbf{Data-plane
pin loopback} bypasses your logic entirely and proves only the adapter, the levels, the wiring and the
pin assignment (not the terminal's baud rate: the adapter sends and receives at the same setting, so a loopback agrees with itself at any rate). \textbf{The control plane} adds the SPI link, the level shifter,
\code{AvrSpi} and the register bank, checked by writing \code{BAUD_DIV} and reading it back.
\textbf{Peripheral loopback} adds the whole datapath, with \code{tx} tied to \code{rx} on the board
so that the byte never leaves the FPGA. \textbf{The real data plane} adds the outside world, and with
it two independently generated baud rates that have to agree. And \textbf{EchoNode} adds the
application on top.

The order is forced, not arbitrary. The peripheral cannot transmit until \code{BAUD_DIV} is written,
and \code{BAUD_DIV} is only reachable over SPI, so the control plane has to be proven before
anything that uses the peripheral. Any ladder that puts a peripheral rung first is really testing two
layers at once and calling it one.

Climbed in order, the first rung that misbehaves names the layer at fault. That is the point of the
method, and it is what an integration test catches that no host test can.

\subsection{What's ahead}
\secref{c10:sec:echo_node} describes \code{app::EchoNode}, the application you write and host-test
before the bench. \secref{c10:sec:exercises} is the bench procedure: assemble the hardware, program
the FPGA and flash the Nano, climb the ladder, run \code{app::EchoNode}, and trace one byte through
every layer the book built.
